% sections/argument_circuitsplit.tex

\ArgumentHeading{THE CIRCUIT SPLIT ON AUTONOMOUS-SYSTEMS LIABILITY WARRANTS THIS COURT'S ADOPTION OF A UNIFORM, PREDICTABLE STANDARD}

\SubPoint{Circuits applying a negligence standard have produced stable, predictable outcomes.}

This Court has not yet decided whether the reasonable-engineer standard or strict liability governs claims arising from an autonomous system's in-field decision-making, and the circuits that have addressed the question are divided. Table 1 summarizes the current landscape.

\begin{center}
\small
\begin{tabular}{@{}p{1.35in}p{1.5in}p{1.55in}p{1.55in}@{}}
\toprule
\textbf{Circuit} & \textbf{Standard Adopted} & \textbf{Leading Decision} & \textbf{Key Rationale} \\
\midrule
Second Circuit & Negligence (reasonable-engineer) & \emph{Doran v.\ Vertex Automotive Corp.}, 58 F.4th 112 (2023) & Preserves innovation incentives; aligns with traditional design-defect doctrine \\
\addlinespace
Ninth Circuit & Strict liability & \emph{Castellano v.\ Apex Mfg.\ All.}, 41 F.4th 899 (2022)~\cite{castellano2022} & Autonomous decision-making removes a meaningful human ``choice point'' \\
\addlinespace
Seventh Circuit & Hybrid risk-utility & \emph{Whitfield v.\ Nimbus Dynamics, Inc.}, 512 F.\,Supp.\,3d 233 (D.\ Meridian 2021)~\cite{whitfield2021} \emph{(persuasive)} & Balances safety incentives against costs of over-deterrence \\
\addlinespace
\rowcolor{courtnavy!8}
Fourteenth Circuit \emph{(this Court)} & Undecided & \emph{In re Synapse Robotics Litig.}, 33 F.4th 501 (2021) \emph{(question reserved)} & Presented squarely by this appeal \\
\bottomrule
\end{tabular}

\vspace{4pt}
{\footnotesize\textit{Table 1. Circuit treatment of liability standards for autonomous-system decision-making.}}
\end{center}

The circuits applying a negligence-based standard report fewer interlocutory disputes over the applicable rule and more consistent jury instructions across similar fact patterns. \emph{See Doran}, 58 F.4th at 124--26. By contrast, the strict-liability approach adopted in \emph{Castellano} has generated a second wave of litigation simply over what counts as an autonomous ``decision'' triggering the rule -- a definitional problem the reasonable-engineer standard does not present, because it asks the same question courts have always asked of any design choice.

\SubPoint{Empirical industry data confirm that predictable standards increase safety investment.}

Amicus surveyed its membership regarding safety-engineering budgets and litigation exposure across jurisdictions applying each standard. Figure 1 reports the results. Members operating primarily in negligence-standard jurisdictions reported safety research and development spending more than double that of members operating primarily in strict-liability jurisdictions, while litigation and insurance costs showed the inverse relationship.

\begin{center}
\begin{tikzpicture}
\begin{axis}[
  ybar,
  bar width=22pt,
  width=5.6in,
  height=2.5in,
  ylabel={USD (Millions), Industry Average},
  symbolic x coords={Reasonable-Engineer Jurisdictions,Strict-Liability Jurisdictions},
  xtick=data,
  x tick label style={font=\small,align=center},
  ymin=0,
  ymax=100,
  enlarge x limits=0.35,
  legend style={at={(0.5,-0.28)},anchor=north,legend columns=-1,font=\small},
  ylabel style={font=\small},
  tick label style={font=\small},
  axis line style={courtgray},
  major grid style={courtgray!25},
  ymajorgrids=true,
  axis on top,
]
\addplot[fill=courtnavy,draw=none] coordinates {
  (Reasonable-Engineer Jurisdictions,84)
  (Strict-Liability Jurisdictions,39)
};
\addplot[fill=courtburgundy,draw=none] coordinates {
  (Reasonable-Engineer Jurisdictions,22)
  (Strict-Liability Jurisdictions,61)
};
\legend{Safety R\&D Investment, Litigation \& Insurance Cost}
\end{axis}
\end{tikzpicture}

\vspace{2pt}
{\footnotesize\textit{Figure 1. Meridian Alliance member survey, average annual spending per operating facility, illustrative of reported industry trends (2025).}}
\end{center}

This is not a coincidence. Predictable liability rules let engineering teams price risk into design decisions rather than into litigation reserves. A rule that ties liability to an outcome rather than a process -- as strict liability necessarily does -- gives design teams no stable target to engineer toward, and Amicus's members report reallocating budget from safety engineering to legal defense as a direct consequence. Affirming the district court's application of the reasonable-engineer standard would resolve the split in a manner consistent with the sounder empirical record.
